% Options for packages loaded elsewhere
% Options for packages loaded elsewhere
\PassOptionsToPackage{unicode}{hyperref}
\PassOptionsToPackage{hyphens}{url}
%
\documentclass[
  ignorenonframetext,
]{beamer}
\newif\ifbibliography
\usepackage{pgfpages}
\setbeamertemplate{caption}[numbered]
\setbeamertemplate{caption label separator}{: }
\setbeamercolor{caption name}{fg=normal text.fg}
\beamertemplatenavigationsymbolsempty
% remove section numbering
\setbeamertemplate{part page}{
  \centering
  \begin{beamercolorbox}[sep=16pt,center]{part title}
    \usebeamerfont{part title}\insertpart\par
  \end{beamercolorbox}
}
\setbeamertemplate{section page}{
  \centering
  \begin{beamercolorbox}[sep=12pt,center]{section title}
    \usebeamerfont{section title}\insertsection\par
  \end{beamercolorbox}
}
\setbeamertemplate{subsection page}{
  \centering
  \begin{beamercolorbox}[sep=8pt,center]{subsection title}
    \usebeamerfont{subsection title}\insertsubsection\par
  \end{beamercolorbox}
}
% Prevent slide breaks in the middle of a paragraph
\widowpenalties 1 10000
\raggedbottom
\AtBeginPart{
  \frame{\partpage}
}
\AtBeginSection{
  \ifbibliography
  \else
    \frame{\sectionpage}
  \fi
}
\AtBeginSubsection{
  \frame{\subsectionpage}
}
\usepackage{iftex}
\ifPDFTeX
  \usepackage[T1]{fontenc}
  \usepackage[utf8]{inputenc}
  \usepackage{textcomp} % provide euro and other symbols
\else % if luatex or xetex
  \usepackage{unicode-math} % this also loads fontspec
  \defaultfontfeatures{Scale=MatchLowercase}
  \defaultfontfeatures[\rmfamily]{Ligatures=TeX,Scale=1}
\fi
\usepackage{lmodern}

\usetheme[noslidenumber]{sintef}
\ifPDFTeX\else
  % xetex/luatex font selection
\fi
% Use upquote if available, for straight quotes in verbatim environments
\IfFileExists{upquote.sty}{\usepackage{upquote}}{}
\IfFileExists{microtype.sty}{% use microtype if available
  \usepackage[]{microtype}
  \UseMicrotypeSet[protrusion]{basicmath} % disable protrusion for tt fonts
}{}
\makeatletter
\@ifundefined{KOMAClassName}{% if non-KOMA class
  \IfFileExists{parskip.sty}{%
    \usepackage{parskip}
  }{% else
    \setlength{\parindent}{0pt}
    \setlength{\parskip}{6pt plus 2pt minus 1pt}}
}{% if KOMA class
  \KOMAoptions{parskip=half}}
\makeatother


\usepackage{longtable,booktabs,array}
\usepackage{calc} % for calculating minipage widths
\usepackage{caption}
% Make caption package work with longtable
\makeatletter
\def\fnum@table{\tablename~\thetable}
\makeatother
\usepackage{graphicx}
\makeatletter
\newsavebox\pandoc@box
\newcommand*\pandocbounded[1]{% scales image to fit in text height/width
  \sbox\pandoc@box{#1}%
  \Gscale@div\@tempa{\textheight}{\dimexpr\ht\pandoc@box+\dp\pandoc@box\relax}%
  \Gscale@div\@tempb{\linewidth}{\wd\pandoc@box}%
  \ifdim\@tempb\p@<\@tempa\p@\let\@tempa\@tempb\fi% select the smaller of both
  \ifdim\@tempa\p@<\p@\scalebox{\@tempa}{\usebox\pandoc@box}%
  \else\usebox{\pandoc@box}%
  \fi%
}
% Set default figure placement to htbp
\def\fps@figure{htbp}
\makeatother


% definitions for citeproc citations
\NewDocumentCommand\citeproctext{}{}
\NewDocumentCommand\citeproc{mm}{%
  \begingroup\def\citeproctext{#2}\cite{#1}\endgroup}
\makeatletter
 % allow citations to break across lines
 \let\@cite@ofmt\@firstofone
 % avoid brackets around text for \cite:
 \def\@biblabel#1{}
 \def\@cite#1#2{{#1\if@tempswa , #2\fi}}
\makeatother
\newlength{\cslhangindent}
\setlength{\cslhangindent}{1.5em}
\newlength{\csllabelwidth}
\setlength{\csllabelwidth}{3em}
\newenvironment{CSLReferences}[2] % #1 hanging-indent, #2 entry-spacing
 {\begin{list}{}{%
  \setlength{\itemindent}{0pt}
  \setlength{\leftmargin}{0pt}
  \setlength{\parsep}{0pt}
  % turn on hanging indent if param 1 is 1
  \ifodd #1
   \setlength{\leftmargin}{\cslhangindent}
   \setlength{\itemindent}{-1\cslhangindent}
  \fi
  % set entry spacing
  \setlength{\itemsep}{#2\baselineskip}}}
 {\end{list}}
\usepackage{calc}
\newcommand{\CSLBlock}[1]{\hfill\break\parbox[t]{\linewidth}{\strut\ignorespaces#1\strut}}
\newcommand{\CSLLeftMargin}[1]{\parbox[t]{\csllabelwidth}{\strut#1\strut}}
\newcommand{\CSLRightInline}[1]{\parbox[t]{\linewidth - \csllabelwidth}{\strut#1\strut}}
\newcommand{\CSLIndent}[1]{\hspace{\cslhangindent}#1}



\setlength{\emergencystretch}{3em} % prevent overfull lines

\providecommand{\tightlist}{%
  \setlength{\itemsep}{0pt}\setlength{\parskip}{0pt}}



 


\makeatletter
\@ifpackageloaded{caption}{}{\usepackage{caption}}
\AtBeginDocument{%
\ifdefined\contentsname
  \renewcommand*\contentsname{Table of contents}
\else
  \newcommand\contentsname{Table of contents}
\fi
\ifdefined\listfigurename
  \renewcommand*\listfigurename{List of Figures}
\else
  \newcommand\listfigurename{List of Figures}
\fi
\ifdefined\listtablename
  \renewcommand*\listtablename{List of Tables}
\else
  \newcommand\listtablename{List of Tables}
\fi
\ifdefined\figurename
  \renewcommand*\figurename{Figure}
\else
  \newcommand\figurename{Figure}
\fi
\ifdefined\tablename
  \renewcommand*\tablename{Table}
\else
  \newcommand\tablename{Table}
\fi
}
\@ifpackageloaded{float}{}{\usepackage{float}}
\floatstyle{ruled}
\@ifundefined{c@chapter}{\newfloat{codelisting}{h}{lop}}{\newfloat{codelisting}{h}{lop}[chapter]}
\floatname{codelisting}{Listing}
\newcommand*\listoflistings{\listof{codelisting}{List of Listings}}
\makeatother
\makeatletter
\makeatother
\makeatletter
\@ifpackageloaded{caption}{}{\usepackage{caption}}
\@ifpackageloaded{subcaption}{}{\usepackage{subcaption}}
\makeatother

\usepackage{bookmark}
\IfFileExists{xurl.sty}{\usepackage{xurl}}{} % add URL line breaks if available
\urlstyle{same}
\hypersetup{
  pdftitle={Cluster Robust Variance Estimation},
  pdfauthor={Edoardo Di Cosimo},
  hidelinks,
  pdfcreator={LaTeX via pandoc}}



\title{Cluster Robust Variance Estimation}
\author{Edoardo Di Cosimo}
\date{}

\begin{document}
\frame{\titlepage}


\begin{frame}{Introduction}
\phantomsection\label{introduction}
\begin{itemize}
\item
  The variance of the sum of 2 random variables \(X_{1}\) and \(X_{2}\)
  is: \[
  \text{Var}(X_{1})+\text{Var}(X_{2})+
  2\cdot
  \underbrace{\text{cov}(X_{1},X_{2})}_{=0\ \text{if } X_1 \perp X_2}
  \]
\item
  The OLS estimator can be written as the true parameter plus a sample
  (weighted) average of individual scores \[
  \hat{\beta}_\text{ols} = \beta + \frac{\sum_{i}X_{i}u_{i}}{\sum _{i}X_{i}^{2}}
  \]
\end{itemize}
\end{frame}

\begin{frame}{Asymptotic}
\phantomsection\label{asymptotic}
\begin{itemize}
\tightlist
\item
  To conduct hypothesis tests we need the asymptotic distribution of our
  estimators
\item
  I.e. we need to understand how the estimator would behave under
  repeated sampling. \[
  \sqrt{ N }\left( \hat{\beta}_\text{ols} - \beta \right) = \frac{\frac{1}{\sqrt{ N }}\sum_{i}X_{i}u_{i}}{\frac{1}{N}\sum _{i}X_{i}^{2}}
  \]
\end{itemize}
\end{frame}

\begin{frame}
\begin{itemize}
\tightlist
\item
  Focus on the numerator \[
  \frac{1}{\sqrt{ N }}\sum_{i}X_{i}u_{i}
  \]
\item
  This is a sum of \(N\) random varianbles
\item
  If they are all independent: \[
  \xrightarrow{d} \mathcal{N}\left( 0,  \frac{1}{N} \sum_{i}\operatorname{Var}(X_{i}u_{i}) \right)
  \]
\end{itemize}
\end{frame}

\begin{frame}{Which is the problem?}
\phantomsection\label{which-is-the-problem}
\begin{itemize}
\tightlist
\item
  Inside a cluster the scores of 2 individuals are not independent
\item
  Usual White Variance estimator is not consistent \[
  \frac{1}{\sqrt{ N }} \sum_{i=1}^N X_{i}u_{i} \xrightarrow{d} \mathcal{N}\left( 0, \frac{1}{N} \sum_{i=1}^N \text{Var}(X_{i}u_{i}) + \frac{2}{N}\sum_{i \ne}^N\sum_{j}^N \text{cov}(X_{i}u_{i},X_{j}u_{j})\right)
  \]
\end{itemize}
\end{frame}

\begin{frame}
\end{frame}

\section{Literature review}\label{literature-review}

\begin{frame}{Why Cluster?}
\phantomsection\label{why-cluster}
\begin{itemize}
\tightlist
\item
  Traditional view: adjust for \textbf{within-cluster error
  correlation}.\\
\item
  Bias in OLS S.E. grows with: strong error correlation, correlated
  regressors, large clusters.\\
\item
  CRVE relies on \textbf{large G} asymptotics.
\end{itemize}
\end{frame}

\begin{frame}{Fixed G}
\phantomsection\label{fixed-g}
\begin{itemize}
\tightlist
\item
  Ibragimov and Müller (2010): inference can use \textbf{cluster-level
  estimates}.\\
\item
  If they are (approximately) independent \& Gaussian a \textbf{t-test}
  with \(t_{G−1}\) is valid even with heterogeneous variances.
\end{itemize}
\end{frame}

\begin{frame}{The Design-Based Argument}
\phantomsection\label{the-design-based-argument}
\begin{itemize}
\tightlist
\item
  Clustering matters because of \textbf{sampling or assignment design},
  not residual correlation.\\
\item
  Residual correlation is \textbf{neither necessary nor sufficient}.\\
\item
  If neither sampling nor treatment assignment is clustered \textbf{do
  not cluster}.\\
\item
  CRVE is often conservative unless effects are homogeneous or clusters
  are rarely observed.
\end{itemize}
\end{frame}

\begin{frame}{What Are Clustered Data?}
\phantomsection\label{what-are-clustered-data}
\begin{itemize}
\tightlist
\item
  Individuals in the same cluster share \textbf{unobserved
  correlation}.\\
\item
  Observations across clusters are identically distributed (mean 0) but
  \textbf{not independent} within clusters.\\
\item
  Knowing one unit's value gives information about another in the same
  group.
\end{itemize}
\end{frame}

\begin{frame}{Representation of Clustered Data}
\phantomsection\label{representation-of-clustered-data}
\begin{itemize}
\tightlist
\item
  Sample with \(N\) individuals grouped in \(G\) clusters:\\
  \(X_{ig} = (X_{11},\dots,X_{N_1 1},\; X_{12},\dots,X_{N_2 2},\dots,X_{N_G G})\)\\
\item
  For \(i,j\) in the same cluster \(g\):\\
  \[
  P(X_i > E[X] \mid X_j > E[X]) > P(X_i > E[X])
  \]
\item
  Cluster \(g\) can be seen as a \textbf{subpopulation}:\\
  \[
  X_g = \{X_i: i \in g\}
  \]
\end{itemize}
\end{frame}

\begin{frame}
\end{frame}

\begin{frame}{Clustered Correlation from Heterogeneous Effects}
\phantomsection\label{clustered-correlation-from-heterogeneous-effects}
\begin{itemize}
\tightlist
\item
  True model:\\
  \[
  Y_{ig} = X_{ig}\beta_g + u_i
  \]
\item
  Estimated model:\\
  \[
  Y_{ig} = X_{ig}\beta_{\text{APE}} + e_{ig}
  \]
\item
  Residual:\\
  \[
  e_{ig} = X_{ig}(\beta_g - \beta_{\text{APE}}) + u_i
  \]
\end{itemize}
\end{frame}

\begin{frame}{Heterogeneity-Induced Correlation}
\phantomsection\label{heterogeneity-induced-correlation}
\begin{itemize}
\tightlist
\item
  Define cluster deviation:\\
  \[
  d_g = \beta_g - \beta_{\text{APE}}
  \]
\item
  Conditional mean of residuals:\\
  \[
  E_g[e_{ig}] = X_{ig} d_g
  \]
\item
  Total variance:\\
  \[
  \sigma^2_{\text{total}}
  = \operatorname{Var}(X_{ig}d_g + u_i)
  \]
\end{itemize}
\end{frame}

\begin{frame}{Within and Between Variance of Residuals}
\phantomsection\label{within-and-between-variance-of-residuals}
\begin{itemize}
\tightlist
\item
  Within-cluster variance:\\
  \[
  \sigma^2_{\text{within}}
  = \operatorname{Var}_g(E(e_{ig}\mid g))
  \]
\item
  Between-cluster variance:\\
  \[
  \sigma^2_{\text{between}}
  = \sigma^2_{\text{total}} - \sigma^2_{\text{within}}
  = \operatorname{Var}(X_{ig}d_g)
  \]
\item
  Hence, \textbf{heterogeneous \(\beta_g\) generates clustered
  residuals}.
\end{itemize}
\end{frame}

\begin{frame}{Assumptions on the DGP}
\phantomsection\label{assumptions-on-the-dgp}
\begin{itemize}
\tightlist
\item
  We must specify how clusters are sampled.\\
\item
  Two frameworks: \textbf{fixed number of clusters} vs \textbf{growing
  number of clusters}.\\
\item
  In both cases: two-stage sampling

  \begin{enumerate}
  \tightlist
  \item
    sample clusters\\
  \item
    sample individuals within clusters.
  \end{enumerate}
\end{itemize}
\end{frame}

\begin{frame}{Fixed Number of Clusters}
\phantomsection\label{fixed-number-of-clusters}
\begin{itemize}
\tightlist
\item
  First stage happens once:

  \begin{itemize}
  \tightlist
  \item
    draw population parameters

    \begin{itemize}
    \tightlist
    \item
      \(\gamma = (\beta_1,\dots,\beta_G)\)\\
    \item
      cluster-specific distributions of \(X_g\) and \(u_g\)\\
    \end{itemize}
  \item
    all fixed across samples.\\
  \end{itemize}
\item
  Second stage repeats across samples:

  \begin{itemize}
  \tightlist
  \item
    random draw of individuals \((X_{ig}, u_{ig})\) from each cluster
    distribution.\\
  \end{itemize}
\item
  \textbf{Only source of randomness} comes from sampling individuals.
\end{itemize}
\end{frame}

\begin{frame}{Increasing Number of Clusters}
\phantomsection\label{increasing-number-of-clusters}
\begin{itemize}
\tightlist
\item
  Both stages are random at every sample:

  \begin{enumerate}
  \tightlist
  \item
    which clusters are drawn\\
  \item
    which individuals are drawn inside each cluster\\
  \end{enumerate}
\item
  Two sources of variability: cluster selection and unit-level sampling.
\end{itemize}
\end{frame}

\begin{frame}{Endogeneity at the Cluster Level}
\phantomsection\label{endogeneity-at-the-cluster-level}
\begin{itemize}
\tightlist
\item
  We check whether\\
  \[
  \mathbb{E}[X_{ig} c_g] \neq 0
  \] for some cluster component \(c_g\).\\
\item
  OLS requires \textbf{exogeneity}: regressors uncorrelated with
  errors.\\
\item
  This may fail \textbf{within} clusters but hold \textbf{on average
  across clusters}.
\end{itemize}
\end{frame}

\begin{frame}{Aggregate vs Within-Cluster Exogeneity}
\phantomsection\label{aggregate-vs-within-cluster-exogeneity}
\begin{itemize}
\tightlist
\item
  Some clusters may show correlation between \(X_{ig}\) and the
  cluster-specific error component (Mundlak, 1978).\\
\item
  If these correlations cancel out across clusters, pooled OLS remains
  \textbf{consistent} even with within-cluster endogeneity.
\end{itemize}
\end{frame}

\begin{frame}{Variance of Cluster-Level Effects}
\phantomsection\label{variance-of-cluster-level-effects}
\begin{itemize}
\tightlist
\item
  Assume a population of cluster-specific effects: \[
  \beta_g \sim \text{ i.i.d. with } 
  E[\beta_g] = \beta_{\text{APE}},\quad
  \operatorname{Var}(\beta_g) = \sigma^2_\beta.
  \]
\item
  We draw \(G\) coefficients: \[
  \boldsymbol{\beta} = (\beta_1,\dots,\beta_G)'.
  \]
\item
  Then \[
  E[\boldsymbol{\beta}] = 
  \begin{bmatrix}
  \beta_{\text{APE}} \\
  \vdots \\
  \beta_{\text{APE}}
  \end{bmatrix},\quad
  \operatorname{Cov}(\boldsymbol{\beta}) =
  \sigma^2_\beta I_G.
  \]
\end{itemize}
\end{frame}

\begin{frame}{Asymptotics in the Number of Clusters}
\phantomsection\label{asymptotics-in-the-number-of-clusters}
\begin{itemize}
\tightlist
\item
  Sample mean of the cluster effects: \[
  \bar{\beta} = \frac{1}{G}\sum_{g=1}^G \beta_g
  \xrightarrow{p} \beta_{\text{APE}}
  \]
\item
  Central Limit Theorem: \[
  \frac{1}{\sqrt{G}} \sum_{g=1}^G (\beta_g - \beta_{\text{APE}})
  \xrightarrow{d} \mathcal{N}(0,\sigma^2_\beta).
  \]
\item
  If all \(\beta_g\) were observed, we could estimate \(\sigma^2_\beta\)
  via: \[
  \hat{\sigma}^2_\beta
  = \frac{G}{G-1} \sum_{g=1}^G (\beta_g - \bar{\beta})^2.
  \]
\end{itemize}
\end{frame}

\begin{frame}{Monte Carlo: Estimating \(\sigma^2_\beta\)}
\phantomsection\label{monte-carlo-estimating-sigma2_beta}
\begin{itemize}
\tightlist
\item
  Monte Carlo design:

  \begin{itemize}
  \tightlist
  \item
    Fix \(G\) and the distribution of \(\beta_g\).
  \item
    Repeat many draws of \((\beta_1,\dots,\beta_G)\).
  \item
    Compare:

    \begin{itemize}
    \tightlist
    \item
      empirical variance of \(\bar{\beta}\) across repetitions;
    \item
      plug-in estimate \(\hat{\sigma}^2_\beta\) from a single draw.
    \end{itemize}
  \end{itemize}
\item
  The resulting plot shows how the \textbf{estimation error}\\
  \(\hat{\sigma}^2_\beta - \operatorname{Var}(\bar{\beta})\) shrinks as
  \(G\) increases.
\end{itemize}
\end{frame}

\begin{frame}{From True \(\beta_g\) to Estimated \(\hat{\beta}_g\)}
\phantomsection\label{from-true-beta_g-to-estimated-hatbeta_g}
\begin{itemize}
\tightlist
\item
  In practice, \(\beta_g\) is not observed; we estimate it from data.\\
\item
  For each cluster \(g\): \[
  \sqrt{N_g}\,(\hat{\beta}_g - \beta_g)
  \xrightarrow{d}
  \mathcal{N}(0, V_g),
  \] where \[
  V_g = (X_g'X_g)^{-1} E_g[X_g' u_g u_g' X_g] (X_g'X_g)^{-1}.
  \]
\item
  Across clusters we thus have \(G\) asymptotically normal estimators,\\
  each centered at its own \(\beta_g\) with cluster-specific variance
  \(V_g\).
\end{itemize}
\end{frame}

\begin{frame}{Interpretation: Signal + Noise}
\phantomsection\label{interpretation-signal-noise}
\begin{itemize}
\tightlist
\item
  Write each estimator as \[
  \hat{\beta}_g = \beta_g + e_g,
  \] with \[
  \sqrt{N_g}\, e_g \xrightarrow{d} \mathcal{N}(0, V_g),\quad E[e_g]=0.
  \]
\item
  The vector \(\hat{\boldsymbol{\beta}}\) can be viewed as \textbf{noisy
  observations}\\
  of the underlying cluster parameters, with noise coming from
  within-cluster sampling.
\end{itemize}
\end{frame}

\begin{frame}{Monte Carlo: Cluster-Wise Estimation}
\phantomsection\label{monte-carlo-cluster-wise-estimation}
\begin{itemize}
\tightlist
\item
  Using a fixed vector of true parameters \(\{\beta_g\}_{g=1}^G\), we:

  \begin{itemize}
  \tightlist
  \item
    build a population with \(G\) clusters;
  \item
    repeatedly resample individuals within each cluster;
  \item
    re-estimate \(\hat{\beta}_g\) in each replication.
  \end{itemize}
\item
  For each cluster \(g\) we obtain a Monte Carlo distribution of
  \(\hat{\beta}_g\):

  \begin{itemize}
  \tightlist
  \item
    histogram of \(\hat{\beta}_g\),
  \item
    vertical line at the true \(\beta_g\).
  \end{itemize}
\item
  This visualizes \(\hat{\beta}_g = \beta_g + e_g\)\\
  and shows how dispersion differs across clusters.
\end{itemize}
\end{frame}

\begin{frame}{Decomposition Around \(\beta_{\text{APE}}\)}
\phantomsection\label{decomposition-around-beta_textape}
\begin{itemize}
\tightlist
\item
  Model the true effects as \[
  \beta_g = \beta_{\text{APE}} + d_g,\quad
  d_g \sim \text{i.i.d. } \mathcal{N}(0,\sigma^2_\beta).
  \]
\item
  Then \[
  \hat{\beta}_g = \beta_{\text{APE}} + d_g + e_g.
  \]
\item
  The cluster mean equals the pooled OLS estimator: \[
  \hat{\beta}_{\text{pols}}
  = \frac{1}{G}\sum_{g=1}^G \hat{\beta}_g
  = \beta_{\text{APE}}
    + \frac{1}{G}\sum_{g=1}^G d_g
    + \frac{1}{G}\sum_{g=1}^G e_g.
  \]
\item
  As \(G \to \infty\), both averages of \(d_g\) and \(e_g\) converge to
  zero in probability.
\end{itemize}
\end{frame}

\begin{frame}{White Variance at the Cluster Level}
\phantomsection\label{white-variance-at-the-cluster-level}
\begin{itemize}
\tightlist
\item
  For each cluster \(g\), we estimate \(V_g\) using the \textbf{White}
  heteroskedasticity-robust estimator.\\
\item
  Monte Carlo check:

  \begin{itemize}
  \tightlist
  \item
    compute the White variance of \(\hat{\beta}_g\);
  \item
    compare it with the Monte Carlo variance of \(\hat{\beta}_g\)
    (across resamples);
  \item
    look at the histogram of ``White variance -- Monte Carlo variance''.
  \end{itemize}
\item
  The small average discrepancy across clusters supports the
  \textbf{consistency of White} for \(V_g\).
\end{itemize}
\end{frame}

\begin{frame}{Pooled OLS and Averaging Cluster Variances}
\phantomsection\label{pooled-ols-and-averaging-cluster-variances}
\begin{itemize}
\tightlist
\item
  White on the pooled regression can be related to averaging
  cluster-level variances:\\
  \[
  \widehat{\operatorname{Var}}(\hat{\beta})
  \;\approx\;
  \frac{1}{G}
  \left(
    \frac{1}{G}\sum_{g=1}^G \widehat{\operatorname{Var}}(\hat{\beta}_g)
  \right),
  \] echoing the Fama--MacBeth idea of averaging across groups.\\
\item
  This interpretation is valid when there is effectively
  \textbf{one-stage sampling}:

  \begin{itemize}
  \tightlist
  \item
    the set of clusters is fixed,
  \item
    randomness comes only from within-cluster sampling.
  \end{itemize}
\end{itemize}
\end{frame}

\begin{frame}{What the Simulations Show}
\phantomsection\label{what-the-simulations-show}
\begin{itemize}
\tightlist
\item
  \begin{enumerate}
  [(i)]
  \tightlist
  \item
    The asymptotic variance of each \(\hat{\beta}_g\) can be
    consistently estimated using\\
    the cluster-level White estimator.\\
  \end{enumerate}
\item
  \begin{enumerate}
  [(i)]
  \setcounter{enumi}{1}
  \tightlist
  \item
    The pooled White variance behaves like an aggregation of the
    cluster-specific White variances.\\
  \end{enumerate}
\item
  \begin{enumerate}
  [(i)]
  \setcounter{enumi}{2}
  \tightlist
  \item
    Once \textbf{resampling of clusters} is introduced, the variance of
    \(\hat{\beta}_{\text{pols}}\)\\
    no longer equals the simple average of
    \(\widehat{\operatorname{Var}}(\hat{\beta}_g)\).
  \end{enumerate}
\end{itemize}
\end{frame}

\begin{frame}{Decomposing Variation in \(\hat{\beta}_g\)}
\phantomsection\label{decomposing-variation-in-hatbeta_g}
\begin{itemize}
\tightlist
\item
  We want to separate two components:

  \begin{itemize}
  \tightlist
  \item
    \textbf{within-cluster estimation error}:
    \(e_g = \hat{\beta}_g - \beta_g\)
  \item
    \textbf{between-cluster heterogeneity}:
    \(d_g = \beta_g - \beta_{\text{APE}}\)
  \end{itemize}
\item
  Two simulation exercises make this decomposition explicit.
\end{itemize}
\end{frame}

\begin{frame}{Exercise 1: Within-Cluster Sampling Error \(e_g\)}
\phantomsection\label{exercise-1-within-cluster-sampling-error-e_g}
\begin{itemize}
\tightlist
\item
  For several cluster sizes \(N_g \in \{10,100,200,300,1000\}\):

  \begin{itemize}
  \tightlist
  \item
    repeatedly resample \(N_g\) individuals per cluster,
  \item
    re-estimate all \(\hat{\beta}_g\),
  \item
    compute estimation errors \(e_g = \hat{\beta}_g - \beta_g\).
  \end{itemize}
\item
  For each \(N_g\), compute \[
  \mathbb{E}[\operatorname{Var}(e_g)]
  \] by averaging variances across clusters and repetitions.
\end{itemize}
\end{frame}

\begin{frame}{Results: Decay of \(e_g\) Variance}
\phantomsection\label{results-decay-of-e_g-variance}
\begin{itemize}
\tightlist
\item
  The variance of \(e_g\) falls rapidly as \(N_g\) increases.\\
\item
  Larger clusters → more precise cluster-specific estimates.\\
\item
  Sampling noise becomes negligible for large \(N_g\).\\
\item
  (Figure: Decay of within-cluster sampling error.)
\end{itemize}
\end{frame}

\begin{frame}{Exercise 2: Between-Cluster Heterogeneity \(d_g\)}
\phantomsection\label{exercise-2-between-cluster-heterogeneity-d_g}
\begin{itemize}
\tightlist
\item
  Repeatedly draw new sets of \(\beta_g\) from the population
  distribution.\\
\item
  Compute deviations: \[
  d_g = \beta_g - \beta_{\text{APE}}.
  \]
\item
  Record the mean deviation \(\overline{d}_g\) for each repetition.\\
\item
  The variance of these means estimates \textbf{heterogeneity-driven
  dispersion}.
\end{itemize}
\end{frame}

\begin{frame}{Results: Nature of \(d_g\)}
\phantomsection\label{results-nature-of-d_g}
\begin{itemize}
\tightlist
\item
  Variation in \(d_g\) reflects \textbf{structural differences}, not
  sampling noise.\\
\item
  This dispersion does not shrink with \(N_g\).\\
\item
  Heterogeneity persists regardless of within-cluster sample size.
\end{itemize}
\end{frame}

\begin{frame}{Varying the Number of Clusters \(G\)}
\phantomsection\label{varying-the-number-of-clusters-g}
\begin{itemize}
\tightlist
\item
  Repeat the heterogeneity experiment for different \(G\).\\
\item
  For each \(G\), draw many sets of \(\beta_g\), compute \(d_g\), and
  record: \[
  \operatorname{Var}(\overline{d}_g).
  \]
\end{itemize}
\end{frame}

\begin{frame}{Results: Shrinkage of
\(\operatorname{Var}(\overline{d}_g)\)}
\phantomsection\label{results-shrinkage-of-operatornamevaroverlined_g}
\begin{itemize}
\tightlist
\item
  As \(G\) increases,\\
  \[
  \operatorname{Var}(\overline{d}_g)
  \] declines.\\
\item
  Averaging over more clusters stabilizes the estimate of the underlying
  mean parameter.\\
\item
  (Figure: variance of mean deviations vs number of clusters.)
\end{itemize}
\end{frame}

\begin{frame}{Monte Carlo Decomposition of
\(\operatorname{Var}(\hat{\beta}_{\text{ols}})\)}
\phantomsection\label{monte-carlo-decomposition-of-operatornamevarhatbeta_textols}
\begin{itemize}
\tightlist
\item
  Goal: assess how well \[
  \operatorname{Var}(\hat{\beta}_{\text{ols}})
  \approx
  \operatorname{Var}(d_g)
  +
  \mathbb{E}\!\left[ \operatorname{Var}(e_g) \right]
  \] holds in finite samples.
\item
  Two components:

  \begin{itemize}
  \tightlist
  \item
    \(d_g = \beta_g - \beta_{\text{APE}}\) (heterogeneity)
  \item
    \(e_g = \hat{\beta}_g - \beta_g\) (sampling error)
  \end{itemize}
\end{itemize}
\end{frame}

\begin{frame}{Simulation Design}
\phantomsection\label{simulation-design}
\begin{itemize}
\tightlist
\item
  Grid of designs:

  \begin{itemize}
  \tightlist
  \item
    cluster sizes \(N_g \in \{20, 50, 100, 200, 300\}\)\\
  \item
    numbers of clusters \(G \in \{5, 10, 20, 50, 100\}\).
  \end{itemize}
\item
  For each \((N_g, G)\) and each repetition:

  \begin{enumerate}
  \tightlist
  \item
    Draw heterogeneous coefficients \(\beta_g\).
  \item
    Generate a population with \(G\) clusters.
  \item
    Sample \(N_g\) individuals from each cluster.
  \item
    Estimate all \(\hat{\beta}_g\) and pooled
    \(\hat{\beta}_{\text{ols}}\).
  \end{enumerate}
\end{itemize}
\end{frame}

\begin{frame}{Recorded Quantities}
\phantomsection\label{recorded-quantities}
\begin{itemize}
\tightlist
\item
  For every repetition and design \((N_g,G)\):

  \begin{itemize}
  \tightlist
  \item
    \(e_g = \hat{\beta}_g - \beta_g\) (estimation errors)
  \item
    \(d_g = \beta_g - \beta_{\text{APE}}\) (heterogeneity terms)
  \item
    \(\hat{\beta}_{\text{ols}}\) (pooled estimator)
  \item
    White and CRVE variance estimates of \(\hat{\beta}_{\text{ols}}\).
  \end{itemize}
\item
  For each \((N_g,G)\) we compute:

  \begin{itemize}
  \tightlist
  \item
    \(\operatorname{Var}(\hat{\beta}_{\text{ols}})\) (empirical
    variance)
  \item
    \(\operatorname{Var}(d_g)\) and
    \(\mathbb{E}[\operatorname{Var}(e_g)]\)
  \item
    relative discrepancy between left- and right-hand side.
  \end{itemize}
\end{itemize}
\end{frame}

\begin{frame}{Performance of the Decomposition}
\phantomsection\label{performance-of-the-decomposition}
\begin{itemize}
\tightlist
\item
  When either \(G\) or \(N_g\) is large:

  \begin{itemize}
  \tightlist
  \item
    \(d_g\) and \(e_g\) are measured very precisely.
  \item
    \(\operatorname{Var}(d_g) + \mathbb{E}[\operatorname{Var}(e_g)]\)\\
    closely matches \(\operatorname{Var}(\hat{\beta}_{\text{ols}})\).
  \end{itemize}
\item
  When both \(G\) and \(N_g\) are small:

  \begin{itemize}
  \tightlist
  \item
    cluster-level estimates are noisy;
  \item
    between-cluster dispersion is overstated;
  \item
    the decomposition is a poorer approximation.
  \end{itemize}
\end{itemize}
\end{frame}

\begin{frame}{Relative Error Across Designs}
\phantomsection\label{relative-error-across-designs}
\begin{itemize}
\tightlist
\item
  The plot of the relative error over \((N_g,G)\) shows:

  \begin{itemize}
  \tightlist
  \item
    largest discrepancies when both \(G\) and \(N_g\) are small;
  \item
    the mismatch rapidly shrinks as we increase either

    \begin{itemize}
    \tightlist
    \item
      the number of clusters, or
    \item
      the cluster size.
    \end{itemize}
  \end{itemize}
\item
  Conclusion: the variance decomposition is \textbf{very accurate} in
  rich designs (many clusters and/or large cluster sizes), but less so
  in sparse settings.
\end{itemize}
\end{frame}

\begin{frame}{CRVE--White Differences as Heterogeneity Measure}
\phantomsection\label{crvewhite-differences-as-heterogeneity-measure}
\begin{itemize}
\tightlist
\item
  We test whether the \textbf{difference} between CRVE and White
  variance estimators\\
  captures the heterogeneity component \[
  \operatorname{Var}(\bar d_g),\quad
  d_g = \beta_g - \beta_{\text{APE}}.
  \]
\item
  For each design \((N_g, G)\):

  \begin{itemize}
  \tightlist
  \item
    compute average CRVE variance,
  \item
    compute average White variance,
  \item
    take their difference,
  \item
    compare it to \(\operatorname{Var}(\bar d_g)\) from the simulations.
  \end{itemize}
\end{itemize}
\end{frame}

\begin{frame}{Simulation-Based Comparison}
\phantomsection\label{simulation-based-comparison}
\begin{itemize}
\tightlist
\item
  For every \((N_g, G)\):

  \begin{itemize}
  \tightlist
  \item
    horizontal axis: \(\operatorname{Var}(\bar d_g)\) (variance of mean
    deviations),
  \item
    vertical axis: \[
    \mathbb{E}\big[\widehat{\operatorname{Var}}_{\text{CRVE}}\big]
    -
    \mathbb{E}\big[\widehat{\operatorname{Var}}_{\text{White}}\big].
    \]
  \end{itemize}
\item
  Each point is a design \((N_g,G)\):

  \begin{itemize}
  \tightlist
  \item
    color intensity encodes \(G\) (lighter = larger \(G\)).
  \end{itemize}
\item
  The 45-degree dashed line is the \textbf{benchmark} of perfect
  equality.
\end{itemize}
\end{frame}

\begin{frame}{Results: Tight Alignment}
\phantomsection\label{results-tight-alignment}
\begin{itemize}
\tightlist
\item
  Points lie very close to the 45-degree line:

  \begin{itemize}
  \tightlist
  \item
    CRVE--White differences almost exactly recover
    \(\operatorname{Var}(\bar d_g)\).
  \end{itemize}
\item
  Designs with \textbf{large \(G\)} (lighter points):

  \begin{itemize}
  \tightlist
  \item
    sit nearly on the identity line,
  \item
    indicate very precise estimation of the heterogeneity component.
  \end{itemize}
\item
  Designs with \textbf{small \(G\)}:

  \begin{itemize}
  \tightlist
  \item
    show slightly more dispersion,
  \item
    but still follow the same linear relationship.
  \end{itemize}
\end{itemize}
\end{frame}

\begin{frame}{Interpretation}
\phantomsection\label{interpretation}
\begin{itemize}
\tightlist
\item
  The gap between CRVE and White is driven by \textbf{between-cluster
  variation},\\
  not by sampling noise.
\item
  Overall, the simulations support: \[
  \mathbb{E}\big[\widehat{\operatorname{Var}}_{\text{CRVE}}\big]
  -
  \mathbb{E}\big[\widehat{\operatorname{Var}}_{\text{White}}\big]
  \approx
  \operatorname{Var}(\bar d_g),
  \] and this relationship strengthens as the number of clusters \(G\)
  grows.
\end{itemize}
\end{frame}

\begin{frame}{Clustered-Robust Variance Estimator (CRVE)}
\phantomsection\label{clustered-robust-variance-estimator-crve}
\begin{itemize}
\tightlist
\item
  Units: \(i=1,\dots,N_g\) within clusters \(g=1,\dots,G\).\\
\item
  Each observation has regressors: \[
  X_{ig} = (x_{1ig},\dots,x_{kig}).
  \]
\item
  Cluster regressor matrix: \[
  X_g =
  \begin{bmatrix}
  X_{1g}\\\vdots\\X_{N_g g}
  \end{bmatrix}
  \ \ (N_g\times k)
  \]
\item
  Full sample: \[
  X =
  \begin{bmatrix}
  X_1\\\vdots\\X_G
  \end{bmatrix}
  \ \ (N\times k),\quad N=\sum_g N_g.
  \]
\end{itemize}
\end{frame}

\begin{frame}{Why Cluster-Robust Inference?}
\phantomsection\label{why-cluster-robust-inference}
\begin{itemize}
\tightlist
\item
  Dependence is allowed \textbf{within} clusters.\\
\item
  Clusters are assumed \textbf{independent across \(g\)}.\\
\item
  Asymptotics rely on \textbf{\(G\to\infty\)}, not on total sample size
  \(N\).
\item
  Following Ibragimov and Müller (2010): consistency comes from
  averaging over many essentially independent groups.
\end{itemize}
\end{frame}

\begin{frame}{Consistency of OLS Under Clustering}
\phantomsection\label{consistency-of-ols-under-clustering}
\begin{itemize}
\tightlist
\item
  OLS: \[
  \hat\beta_{\text{ols}} = (X'X)^{-1} X'y
  = \left(\sum_g X_g'X_g\right)^{-1}\!\left(\sum_g X_g'y_g\right)
  \]
\item
  Error representation: \[
  \hat\beta - \beta
  =
  \left(\tfrac1G\sum_g X_g'X_g\right)^{-1}
  \left(\tfrac1G\sum_g X_g'u_g\right)
  \]
\item
  If\\
  \[
  E[X_g' u_g] = 0
  \] then by WLLN: \[
  \tfrac1G\sum_g X_g'u_g \xrightarrow{p}0
  \] so \(\hat\beta_{\text{ols}}\) is consistent.
\end{itemize}
\end{frame}

\begin{frame}{Asymptotic Distribution of OLS}
\phantomsection\label{asymptotic-distribution-of-ols}
\begin{itemize}
\tightlist
\item
  Scaling: \[
  \sqrt{G}(\hat\beta-\beta)
  =
  \left(\tfrac1G\sum_g X_g'X_g\right)^{-1}
  \left(\tfrac{1}{\sqrt{G}}\sum_g X_g'u_g\right)
  \]
\item
  First term \(\to A^{-1}\).\\
\item
  Second term is asymptotically normal: \[
  \tfrac1{\sqrt G}\sum_g X_g'u_g
  \overset{a}{\sim}
  \mathcal{N}\bigl(0,\,E[X_g' u_g u_g' X_g]\bigr)
  \]
\item
  Therefore: \[
  \sqrt{G}(\hat\beta-\beta)
  \xrightarrow{d}
  \mathcal{N}\!\left(0,\; A^{-1} E[X_g'u_g u_g'X_g] A^{-1}\right)
  \]
\end{itemize}
\end{frame}

\begin{frame}{CRVE: General Formula}
\phantomsection\label{crve-general-formula}
\begin{itemize}
\tightlist
\item
  Cluster-robust variance matrix: \[
  \hat V_{\text{clu}}(\hat\beta)
  =
  (X'X)^{-1}
  \left(\sum_{g=1}^G X_g'\hat u_g \hat u_g' X_g\right)
  (X'X)^{-1}
  \]
\item
  Key object: \[
  s_g = X_g'\hat u_g
  \] so that\\
  \[
  s_g s_g' = X_g' \hat u_g \hat u_g' X_g.
  \]
\end{itemize}
\end{frame}

\begin{frame}{Cluster-Level Contribution}
\phantomsection\label{cluster-level-contribution}
\begin{itemize}
\tightlist
\item
  For a cluster: \[
  s_g s_g'
  =
  \bigl[\sum_{i,j\in g} x_{aig}x_{bjg}\,\hat u_{ig}\hat u_{jg}\bigr]_{a,b=1}^k
  \]
\item
  Summing over \(g\) gives the full middle matrix: \[
  \hat B = \sum_g s_g s_g'.
  \]
\end{itemize}
\end{frame}

\begin{frame}{Convergence of CRVE}
\phantomsection\label{convergence-of-crve}
\begin{itemize}
\tightlist
\item
  Start with: \[
  \tfrac1G \sum_g X_g'\hat u_g \hat u_g' X_g
  \]
\item
  Replace \(\hat u_g = y_g - X_g\hat\beta\) and expand.\\
\item
  Terms involving \(\beta-\hat\beta\) vanish as \(\hat\beta\to\beta\).\\
\item
  Thus: \[
  \tfrac1G\sum_g X_g'\hat u_g\hat u_g'X_g
  \xrightarrow{p}
  \tfrac1G\sum_g X_g'u_g u_g' X_g
  \xrightarrow{p}
  E[X_g'u_g u_g'X_g]
  \]
\item
  Hence CRVE converges to the asymptotic variance of OLS.
\end{itemize}
\end{frame}

\begin{frame}{CRVE vs White: Selector Matrix View}
\phantomsection\label{crve-vs-white-selector-matrix-view}
\begin{itemize}
\tightlist
\item
  Start from: \[
  \hat S\hat S' = X'\hat u \hat u' X.
  \]
\item
  \textbf{White} keeps only diagonal elements (\(i=j\)): \[
  B_W = \operatorname{diag}(\hat u_1^2,\dots,\hat u_N^2).
  \]
\item
  \textbf{CRVE} keeps all \((i,j)\) pairs \textbf{within the same
  cluster}, zeroing cross-cluster terms.
\end{itemize}
\end{frame}

\begin{frame}{Example: One Cluster of Three Observations}
\phantomsection\label{example-one-cluster-of-three-observations}
\begin{itemize}
\tightlist
\item
  \textbf{White}: \[
  B_W =
  \begin{bmatrix}
  \hat u_1^2 & 0 & 0\\
  0 & \hat u_2^2 & 0\\
  0 & 0 & \hat u_3^2
  \end{bmatrix}
  \]
\item
  \textbf{CRVE}: \[
  B_C =
  \begin{bmatrix}
  \hat u_1^2 & \hat u_1\hat u_2 & \hat u_1\hat u_3\\
  \hat u_1\hat u_2 & \hat u_2^2 & \hat u_2\hat u_3\\
  \hat u_1\hat u_3 & \hat u_2\hat u_3 & \hat u_3^2
  \end{bmatrix}
  \]
\end{itemize}
\end{frame}

\begin{frame}{Contribution to the Variance Estimator}
\phantomsection\label{contribution-to-the-variance-estimator}
\begin{itemize}
\tightlist
\item
  After multiplying by \(X_g\): \[
  X_g' B_W X_g
  = \sum_{i\in g} x_i^2\hat u_i^2
  \]
\item
  CRVE adds all extra within-cluster covariances: \[
  X_g' B_C X_g
  =
  X_g' B_W X_g 
  +
  2\!\sum_{i<j\in g} x_i x_j\,\hat u_i\hat u_j.
  \]
\end{itemize}
\end{frame}

\begin{frame}{Key Difference}
\phantomsection\label{key-difference}
\begin{itemize}
\tightlist
\item
  \textbf{White} ignores intra-cluster dependence.\\
\item
  \textbf{CRVE} adds all within-cluster cross-products: \[
  x_i x_j\,\hat u_i\hat u_j.
  \]
\item
  These extra terms adjust the variance for correlated errors inside
  clusters.\\
\item
  This is the exact mechanism through which CRVE corrects inference when
  within-cluster dependence is present.
\end{itemize}
\end{frame}

\begin{frame}{Main Findings}
\phantomsection\label{main-findings}
\begin{itemize}
\tightlist
\item
  The thesis combines theory and simulation to study cluster-robust
  variance estimation.\\
\item
  Simulations clarify when clustering matters, how pooled OLS behaves
  under heterogeneous effects, and when White or CRVE should be used.\\
\item
  Results offer a unified view of clustered sampling, estimator
  consistency, and variance behavior.
\end{itemize}
\end{frame}

\begin{frame}{Consistency of Cluster-Specific Estimators}
\phantomsection\label{consistency-of-cluster-specific-estimators}
\begin{itemize}
\tightlist
\item
  With no cluster-level endogeneity:

  \begin{itemize}
  \tightlist
  \item
    each \(\hat{\beta}_g\) is consistent for its \(\beta_g\);\\
  \item
    including cluster intercepts absorbs within-cluster dependence;\\
  \item
    White provides accurate variance estimates at the cluster level.
  \end{itemize}
\item
  The average partial effect\\
  \[
  \beta_{\text{APE}}
  \] is consistently estimated by averaging \(\hat{\beta}_g\) when all
  clusters are observed.
\end{itemize}
\end{frame}

\begin{frame}{Cluster-Level Endogeneity and Pooled OLS}
\phantomsection\label{cluster-level-endogeneity-and-pooled-ols}
\begin{itemize}
\tightlist
\item
  Regressors may correlate with shocks \textbf{within} some clusters.\\
\item
  Exogeneity may still hold \textbf{on average across clusters},
  preserving consistency of pooled OLS.
\item
  Variance of \(\hat{\beta}_{\text{ols}}\) reflects:

  \begin{itemize}
  \tightlist
  \item
    between-cluster heterogeneity: \(d_g\),\\
  \item
    within-cluster sampling error: \(e_g\).
  \end{itemize}
\end{itemize}
\end{frame}

\begin{frame}{Variance Decomposition}
\phantomsection\label{variance-decomposition}
\begin{itemize}
\tightlist
\item
  Supported strongly by simulations: \[
  \mathrm{Var}(\hat{\beta}_{\text{ols}})
  \approx
  \mathrm{Var}(d_g) + \mathbb{E}[\mathrm{Var}(e_g)].
  \]
\item
  Works extremely well when:

  \begin{itemize}
  \tightlist
  \item
    \(G\) is large, \textbf{or}
  \item
    \(N_g\) is large.
  \end{itemize}
\item
  Breaks down only in small designs (few clusters \& small cluster
  sizes).
\end{itemize}
\end{frame}

\begin{frame}{White vs CRVE}
\phantomsection\label{white-vs-crve}
\begin{itemize}
\tightlist
\item
  When there is no between-cluster heterogeneity:

  \begin{itemize}
  \tightlist
  \item
    CRVE = White.
  \end{itemize}
\item
  With heterogeneity:

  \begin{itemize}
  \tightlist
  \item
    CRVE exceeds White by ≈ \(\operatorname{Var}(d_g)\).
  \item
    The CRVE--White gap is informative about the magnitude of \(d_g\).
  \end{itemize}
\item
  But:

  \begin{itemize}
  \tightlist
  \item
    CRVE is \textbf{not} a generic robustness tool.
  \item
    If clustering is not implied by sampling/assignment, CRVE becomes
    conservative and misestimates the true variance.
  \end{itemize}
\end{itemize}
\end{frame}

\begin{frame}{Final Takeaways}
\phantomsection\label{final-takeaways}
\begin{enumerate}
\tightlist
\item
  \textbf{Heterogeneous treatment effects} do not threaten consistency
  of pooled OLS under population exogeneity.\\
\item
  \textbf{Cluster variance decomposes naturally} into \(d_g\)
  (heterogeneity) and \(e_g\) (sampling noise), with simulations closely
  matching theory.\\
\item
  \textbf{CRVE is valid only when justified by the DGP}---not simply by
  observing correlated residuals.
\end{enumerate}

\begin{itemize}
\tightlist
\item
  Theory + simulation together show when cluster-robust inference is
  appropriate, informative, and reliable.
\end{itemize}
\end{frame}

\begin{frame}
\phantomsection\label{refs}
\begin{CSLReferences}{1}{0}
\bibitem[\citeproctext]{ref-ibragimov2010t}
Ibragimov, Rustam, and Ulrich K Müller. 2010. {``T-Statistic Based
Correlation and Heterogeneity Robust Inference.''} \emph{Journal of
Business \& Economic Statistics} 28 (4): 453--68.

\end{CSLReferences}
\end{frame}




\end{document}
